\subsection{LIE POLYNOMIALS AND SUBSTITUTIONS}

Let I be a set. Let \(\mathrm{T}_i\) denote the canonical image of the element \(i\) of I in \(\mathbf{L}(\mathbf{I})\) (which is also sometimes denoted by \(\mathrm{L}((\mathrm{T}_i)_{i\in \mathbf{I}})\) ); the elements of \(\mathrm{L}(\mathbf{I})\) are called Lie polynomials in the indeterminates \((\mathrm{T}_i)_{i\in \mathbf{I}}\) .

Let g be a Lie algebra. If \(\boldsymbol{t} = (t_{i})_{i \in I}\) is a family of elements of g, let \(f_{\boldsymbol{t}}\) denote the homomorphism of L(I) into g such that \(f_{\boldsymbol{t}}(\mathrm{T}_{\boldsymbol{i}}) = t_{i}\) for \(i \in I\) (no. 2, Proposition 1). The image under \(f_{\boldsymbol{t}}\) of the element P of L(I) is denoted by \(\mathrm{P}((t_{i})_{i \in I})\) . In particular, \(\mathrm{P}((\mathrm{T}_{\boldsymbol{i}})_{i \in I}) = \mathrm{P}\) ; the above element \(\mathrm{P}((t_{i})_{i \in I})\) is sometimes called the element of g obtained by substituting the \(t_{i}\) for the \(T_{i}\) in the Lie polynomial \(\mathrm{P}((\mathrm{T}_{\boldsymbol{i}})_{i \in I})\) .

Let \(\sigma: \mathfrak{g} \to \mathfrak{g}'\) be a Lie algebra homomorphism. For every family \(t = (t_i)_{i \in I}\) of elements of \(\mathfrak{g}\) and all \(P \in L(I)\) ,

\[
\sigma (\mathrm{P} ((t _ {i}) _ {i \in \mathrm{I}})) = \mathrm{P} ((\sigma (t _ {i})) _ {i \in \mathrm{I}}),\tag{4}
\]

for \(\sigma^{\circ}f_{t}\) maps \(\mathrm{T}_i\) to \(\sigma(t_i)\) for \(i\in I\) .

Let \((\mathbf{Q}_j)_{j\in J}\) be a family of elements of \(\mathbf{L}(\mathbf{I})\) and let \(\mathbf{P}\in \mathbf{L}(\mathbf{J})\) . By substituting the \(Q_{j}\) for the \(T_{j}\) in \(P\) , we obtain a Lie polynomial \(R = P((Q_{j})_{j \in J}) \in L(I)\) . Then

\[
\mathrm{R} ((t _ {i}) _ {i \in \mathrm{I}}) = \mathrm{P} ((Q _ {j} ((t _ {i}) _ {i \in \mathrm{I}})) _ {j \in J})\tag{5}
\]

for every family \(\pmb{t} = (t_i)_{i\in \mathbf{I}}\) of elements of a Lie algebra \(\mathfrak{g}\) , as is seen by operating by the homomorphism \(f_{t}\) on the equation \(\mathrm{R} = \mathrm{P}((\mathrm{Q}_j)_{j\in \mathbf{J}})\) and using (4).

Let \(\mathfrak{g}\) be a Lie algebra, I a finite set and \(P \in L(I)\) . Suppose that \(\mathfrak{g}\) is a free K-module. The mapping

\[
\tilde {\mathrm{P}} \colon \mathfrak {g} ^ {\mathrm{I}} \to \mathfrak {g}
\]

defined by \(\tilde{\mathrm{P}}((t_i)_{i\in \mathbf{I}}) = \mathrm{P}((t_i)_{i\in \mathbf{I}})\) is then polynomial.† For the set F of mappings of \(\mathfrak{g}^{\mathbf{I}}\) into \(\mathfrak{g}\) is a Lie algebra with the bracket defined by

\[
[ \phi , \psi ] (\boldsymbol {t}) = [ \phi (\boldsymbol {t}), \psi (\boldsymbol {t}) ];\tag{6}
\]

the set \(\mathbf{F}'\) of polynomial mappings of \(\mathfrak{g}^{\intercal}\) into \(\mathfrak{g}\) is a Lie subalgebra of \(\mathbf{F}\) by the bilinearity of the bracket. Our assertion then follows from the fact that the mapping \(\mathrm{P} \mapsto \tilde{\mathrm{P}}\) is a Lie algebra homomorphism and \(\tilde{\mathrm{T}}_i = \mathrm{pr}_i \in \mathrm{F}'\) for all \(i\) .

